\section{Signaling \& Cheap Talk}

\subsection{Signaling Games}
Sender knows type \(\theta \in \Theta\) (prior \(p\)), chooses costly message \(m\); receiver observes \(m\), updates beliefs \(\mu(\theta|m)\), takes action \(a\).

\textbf{PBE requirements:}
\begin{enumerate}[nosep]
    \item Sender: \(m^*(\theta)\) optimal given receiver strategy \(a^*(m)\)
    \item Receiver: \(a^*(m)\) optimal given beliefs \(\mu(\cdot|m)\)
    \item Beliefs: Bayes' rule on the equilibrium path; \textbf{unrestricted off path}
\end{enumerate}

\textbf{Equilibrium taxonomy:}
\begin{itemize}[nosep]
    \item \textbf{Separating:} \(m^*\) injective---types revealed. Requires \textbf{single crossing}: signal cost falls with type, so low types won't mimic
    \item \textbf{Pooling:} all types send the same \(m\)---beliefs at the pooled message equal the prior; sustained by pessimistic off-path beliefs
    \item \textbf{Semi-separating:} some types mix
\end{itemize}

\textbf{Spence education model:} productivity \(\theta_H > \theta_L\), signal cost \(c(s,\theta) = s/\theta\), competitive wages \(w(s) = \mathbb{E}[\theta|s]\). Separating equilibrium: \(s_L = 0\), \(s_H = s^*\) with
\begin{gather*}
\theta_H - s^*/\theta_L \leq \theta_L \quad \text{(L won't mimic)} \\
\theta_H - s^*/\theta_H \geq \theta_L \quad \text{(H prefers signaling)}
\end{gather*}
Education can be pure waste and still be individually rational.

\textbf{Refinement (Intuitive Criterion, Cho--Kreps):} prune off-path beliefs that put weight on types who could not gain from the deviation under \textit{any} receiver response. Selects the \textbf{least-cost separating} (Riley) equilibrium in Spence; kills pooling equilibria supported by unreasonable beliefs.

\subsection{Cheap Talk (Crawford--Sobel)}
Messages are \textbf{costless and nonbinding}: information transmission relies purely on aligned interests, not signal cost.

\textbf{Setup:} state \(\theta \sim U[0,1]\) known to sender; receiver picks \(a\); payoffs \(-(a-\theta)^2\) (receiver), \(-(a-\theta-b)^2\) (sender), bias \(b > 0\).

\textbf{Results:}
\begin{itemize}[nosep]
    \item \textbf{Babbling equilibrium always exists}: messages ignored, no information sent---cheap talk games never have a unique equilibrium
    \item All equilibria are \textbf{partitional}: sender reveals only which interval \([\theta_i, \theta_{i+1})\) contains \(\theta\)
    \item Partition steps satisfy \(\theta_{i+1} - \theta_i = \theta_i - \theta_{i-1} + 4b\): intervals lengthen (bias compounds upward)
    \item Max number of intervals \(N(b)\) is finite, decreasing in \(b\); as \(b \to 0\), \(N \to \infty\) (full revelation in the limit)
    \item No communication possible when \(b \geq 1/4\)
\end{itemize}

\textbf{Contrast with signaling:} cost separates types in Spence; \textit{preference alignment} separates states here. Committing to coarse language is what makes partial honesty credible.

\textit{Application:} expert advice (analysts, doctors, advisors)---the more the expert's incentives diverge, the vaguer the credible advice.
